
The various kinematic distributions and fiducial regions used to fit $m_W$ in the ATLAS and
Tevatron experiments provides some sensitivity to PDF predictions.  Other $W$- and $Z$-boson 
measurements from the LHC and the 
Tevatron provide more significant PDF constraints and are used in the determination of the 
PDF sets.  This section compares the compatibility of these other measurements with the
various PDF sets.  Some sets have low compatibility and are not favoured for an $m_W$
combination.

The $W$-boson rapidity ($y_W$) distribution affects the $m_W$ measurement through the
\pTl~distribution: more central $W$ bosons can have more forward-decaying leptons within the 
detector acceptance, lowering the mean observed \pTl.  Measurements that probe PDF parameters 
describing $y_W$ include the $Z$ boson rapidity $y_Z$ and the asymmetries in the rapidity 
distribution between positive and negative $W$ bosons ($A_W$), or similarly the positive and 
negative charged leptons from their decay ($A_\ell$).  These measurements are considered in this 
compatibility study, and are shown in Table~\ref{tab:datasets}.  

The comparison between data and predictions is performed with the xFitter~\cite{Bertone:2017tig} 
framework.  A $\chi^2$ measure is constructed including all experimental uncertainties and their 
correlations, as well as the PDF uncertainties.
Theory predictions
are calculated at NNLO in QCD and corrected to NLO electroweak predictions using multiplicative $k$-factors
in each measurement bin.  PDF uncertainties are computed at NLO in QCD using Applgrids~\cite{Carli:2010rw}
with calculations from MCFM-6.8~\cite{Campbell:1999ah}.
The results for various PDF sets are shown in Table~\ref{tab:pdfchi2}.

Most of the Drell-Yan measurements have good $\chi^2$ values for all PDFs.  The most significant 
outlier is the D0 $W\rightarrow e \nu$ lepton asymmetry measurement, for which the CT18 set has 
the lowest $\chi^2$ primarily due to its larger uncertainties.  These larger uncertainties also
reduce the correlated $\chi^2$, which represents the contribution from correlated
uncertainties~\cite{alekhin:2015herafitter}.
The correlated $\chi^2$ reduces from 251 to 43 after including
PDF uncertainties in the CT18 set; the corresponding reduction for the NNPDF3.1 set is 110 to 76.
The overall probability of consistency of the combined datasets is 1.5\% for the CT18 set, and is
much lower for the other sets.
Among the studied PDF sets CT18 is therefore considered to give the most accurate estimate of the
68\% C.L. interval for combined $W$- and $Z$-boson measurements.

\begin{table}[!tbp]
  \begin{center}
    \begin{tabular}{cccllc}
\hline \hline
Exp.                       & Obs.     & Decay    & $\sqrt{s}$ & Lum.           & bins \\
\hline
CDF~\cite{CDF:2009cjw}     & $A_W$    & $e\nu$   & 1.96~TeV   & 1~fb$^{-1}$    & 13 \\ %& arXiv:0901.2169 \\
CDF~\cite{CDF:2010vek}     & $y_Z$    & $ee$     & 1.96~TeV   & 2.1~fb$^{-1}$  & 28 \\ %& arXiv:0908.3914 \\
\hline
D0~\cite{D0:2007djv}       & $y_Z$    & $ee$     & 1.96~TeV   & 0.4~fb$^{-1}$  & 28 \\ %& hep-ex/0702025 \\
D0~\cite{D0:2013xqc}       & $A_\ell$ & $\mu\nu$ & 1.96~TeV   & 7.3~fb$^{-1}$  & 12 \\ %& arXiv:1309.2591 \\
D0~\cite{D0:2014kma}       & $A_\ell$ & $e\nu$   & 1.96~TeV   & 9.7~fb$^{-1}$  & 13 \\ % & arXiv:1412.2862 \\
\hline
ATLAS~\cite{ATLAS:2016nqi} & $Z$,$W$ & $\ell\ell$,$\ell\nu$ & 7~TeV & 4.7~fb$^{-1}$ & 61 \\ %& arXiv:1612.03016 \\
\hline \hline
\end{tabular}
    \caption{Drell-Yan measurements used for the PDF compatibility study.
    } 
      \label{tab:datasets} 
\end{center}
\end{table}

\begin{table*}[!tbp]
  \begin{center}
  \begin{tabular*}{\linewidth}{l @{\extracolsep{\fill}}lllllll}
\hline \hline 
  Measurement                & NNPDF3.1   & NNPDF4.0 & MMHT14   & MSHT20  & CT14      & CT18      & ABMP16  \\
\hline 
%\midrule
  CDF $y_Z$                  & 24 / 28   & 28 / 28   & 30 / 28  & 32 / 28 & 29 / 28   & 27 / 28   & 31 / 28  \\
  CDF $A_W$                  & 11 / 13   & 14 / 13   & 12 / 13  & 28 / 13 & 12 / 13   & 11 / 13   & 21 / 13  \\
  D0  $y_Z$                  & 22 / 28   & 23 / 28   & 23 / 28  & 24 / 28 & 22 / 28   & 22 / 28   & 22 / 28  \\
  D0 $W \to e\nu$ $A_\ell$   & 22 / 13   & 23 / 13   & 52 / 13  & 42 / 13 & 21 / 13   & 19 / 13   & 26 / 13  \\
  D0 $W \to \mu\nu$ $A_\ell$ & 12 / 10   & 12 / 10   & 11 / 10  & 11 / 10 & 11 / 10   & 12 / 10   & 11 / 10  \\
  ATLAS peak CC $y_Z$        & 13 / 12   & 13 / 12   & 58 / 12  & 17 / 12 & 12 / 12   & 11 / 12   & 18 / 12  \\
  ATLAS $W^{-}$ $y_\ell$        & 12 / 11   & 12 / 11   & 33 / 11  & 16 / 11 & 13 / 11   & 10 / 11   & 14 / 11  \\
  ATLAS $W^{+}$ $y_\ell$        &  9 / 11   &  9 / 11   & 15 / 11  & 12 / 11 & 9 / 11    &  9 / 11   & 10 / 11  \\
  Correlated $\chi^2$        & 75        & 62        & 210      & 88      & 81        & 41        & 83       \\
\hline
  Total $\chi^2$ / d.o.f.    & 200 / 126 & 196 / 126 & 444 / 126 & 270 / 126 & 210 / 126 & 162 / 126  & 236 / 126  \\
\hline
  p($\chi^2,n$)        & 0.003\%   & 0.007\%   & $<10^{-10}$  & $<10^{-10}$ & $0.0004\%$ & 1.5\% & $10^{-8}$   \\
\hline \hline 
    \end{tabular*}
    \caption{$\chi^2$ per degree of freedom for the Tevatron $Z$-rapidity and $W$- and $l$-asymmetry
    measurements at $\sqrt{s}=1.96$~TeV, and the LHC $Z$-rapidity and $W$ lepton-rapidity measurements at $\sqrt{s}=7$~TeV.  
    The total $\chi^2$ is the sum of those quoted for individual measurements along with a separate contribution for
    correlated uncertainties, where the latter is extracted using a nuisance parameter representation of the
    $\chi^2$~\cite{alekhin:2015herafitter}.  The CT14 and CT18 PDF uncertainties correspond to 68\% coverage, obtained by
    rescaling the eigenvectors by a factor of 1/1.645.  The probability of obtaining a total $\chi^2$ at least as high
    as that observed is labelled p($\chi^2,n$).
}
    \label{tab:pdfchi2}
  \end{center}
\end{table*}

